\setchapterpreamble[u]{\margintoc}
\chapter{T-SQL: Manipulação de Dados (DML)}
\labch{cet615-dml}

Este capítulo cobre a primeira metade da segunda
aula~\cite{BaiocoCET615}: inserir, alterar e excluir dados; consultar
com filtros, agrupamentos, junções e subconsultas; e agrupar operações em
transações. Os exemplos de consulta da aula usam o banco PUBS, de
exemplo da Microsoft. Aqui os exemplos usam o banco de biblioteca do
\refch{cet615-ddl}.

\section{Inserir, alterar e excluir}

\subsection{\texttt{insert}}

\begin{lstlisting}
-- sem lista de colunas: todos os valores, na ordem da tabela
insert into leitor values (1, 'Ana Lima', null, 'Campinas')

-- com lista de colunas: as omitidas recebem nulo ou o default
insert into leitor (codleitor, nome) values (2, 'Bruno Reis')

-- várias linhas no mesmo comando
insert into leitor (codleitor, nome)
values (3, 'Carla Dias'), (4, 'Davi Rocha')

-- resultado de uma consulta
insert into leitor_vip (codleitor, nome)
select codleitor, nome from leitor where cidade = 'Campinas'
\end{lstlisting}

Texto e data vão entre aspas; número, não. \texttt{null} é um valor
especial que significa ``não informado'': não é branco nem zero. Gravar
\texttt{''} ou \texttt{0} grava exatamente isso. No \texttt{insert} sem
lista de colunas, o nulo precisa ser escrito na posição da coluna.

Duas observações da aula\sidecite{BaiocoCET615Aulas}: o \texttt{insert}
com várias linhas tem o mesmo resultado de vários \texttt{insert}, mas
roda como uma única transação; e o \texttt{insert ... select} serve, por
exemplo, para montar uma tabela desnormalizada para consulta rápida ou
para uma ferramenta de BI.

\subsection{\texttt{update} e \texttt{delete}}

\begin{lstlisting}
update livro set titulo = 'Dom Casmurro' where codlivro = 10
update exemplar set estado = 'U', codlivro = 11 where codexemplar = 5
update livro set preco = preco * 1.1        -- todas as linhas!

delete from emprestimo where codemprestimo = 100
delete emprestimo                            -- o from é opcional
\end{lstlisting}

O \texttt{where} é ``quase obrigatório'': sem ele, o comando atinge todas
as linhas. Filtrar pela chave garante uma linha só; filtrar por outra
coluna pode atingir várias. Um \texttt{update} ou \texttt{delete} que não
encontra nada não dá erro: só informa ``0 linhas afetadas''. Esse
comportamento volta nas transações (\refsec{cet615-transacoes-sqlserver}).

\subsection{\texttt{truncate table}}

Para esvaziar uma tabela grande, \texttt{delete} sem \texttt{where} pode
estourar o log de transações, que é bem menor que o banco. O comando
indicado é \texttt{truncate table emprestimo}.

\subsection*{Complemento: o que o \texttt{truncate} registra}

Em aula, a professora disse que o \texttt{truncate} não usa o log; o
slide diz que ele usa menos log. A documentação\sidecite{MSTruncate}
confirma o slide: o \texttt{delete} grava no log uma entrada por linha,
enquanto o \texttt{truncate} desaloca as páginas da tabela e registra só
essas desalocações. Outras diferenças documentadas:

\begin{itemize}
	\item o \texttt{truncate} pode ser desfeito com \texttt{rollback}
	dentro de uma transação;
	\item não funciona em tabela referenciada por chave estrangeira;
	\item não dispara gatilhos, porque não registra as linhas
	excluídas;
	\item reinicia a contagem de colunas \texttt{identity}.
\end{itemize}

\section{Consultas}

\subsection{Estrutura do \texttt{select}}

\begin{lstlisting}
select [distinct] colunas      -- o que mostrar
from tabelas                   -- de onde
[where condição]               -- quais linhas
[group by colunas]             -- como agrupar
[having condição]              -- quais grupos
[order by colunas [asc|desc]]  -- em que ordem
\end{lstlisting}

Antes de consultar, escolha o banco com \texttt{use}. A professora
lembrou que conhecer o esquema (chaves e relacionamentos) é o primeiro
passo para escrever consultas, principalmente com junção.

\subsection{SQL e álgebra relacional}

Toda consulta SQL é traduzida pelo SGBD em operações de álgebra
relacional. Na compilação, o SGBD considera árvores de álgebra
equivalentes e escolhe a mais barata; a otimização não depende só de
índices e estatísticas. A correspondência mostrada em aula:

\begin{description}
	\item[Lista do \texttt{select} = projeção ($\pi$):] corte vertical,
	escolhe colunas.
	\item[\texttt{where} = seleção ($\sigma$):] corte horizontal, escolhe
	linhas.
	\item[Junção no \texttt{from}:] produto cartesiano seguido de uma
	condição, mas executado por um algoritmo próprio, mais eficiente.
\end{description}

Na álgebra, o resultado de uma consulta é uma relação e, portanto, não
tem tuplas repetidas. Os SGBDs não eliminam repetições por padrão (seria
caro), por isso o \texttt{distinct} precisa ser pedido.
\textcite[cap. 6]{Elmasri2011} cobre a álgebra relacional, que a
professora recomendou revisar.

\subsection{Filtros}

\begin{lstlisting}
select titulo from livro where preco > 50
select titulo from livro where preco between 20 and 50
select nome from leitor where cidade in ('Campinas', 'Jundiaí')
select nome from leitor where cidade not in ('Campinas')
select nome from leitor where nome like 'A%'      -- começa com A
select nome from leitor where nome like '_na%'    -- 2.a e 3.a letras "na"
select nome from leitor where nome like '[A-C]%'  -- começa com A, B ou C
\end{lstlisting}

\texttt{between}, \texttt{in} e \texttt{like} com colchetes são
atalhos: tudo pode ser escrito com \texttt{>=}, \texttt{<=}, \texttt{=}
e \texttt{or}, só que mais longo. No \texttt{like}, \texttt{\%} vale
qualquer sequência (inclusive vazia) e \texttt{\_} vale exatamente um
caractere naquela posição, o que fixa o tamanho do texto.

A professora frisou a \textbf{precedência}: operadores relacionais vêm
antes dos lógicos, e \texttt{and} vem antes de \texttt{or}.

\begin{lstlisting}
-- errado: o filtro de preço só vale para 'romance'
select titulo from livro
where genero = 'poesia' or genero = 'romance' and preco > 50

-- certo
select titulo from livro
where (genero = 'poesia' or genero = 'romance') and preco > 50
\end{lstlisting}

Parênteses a mais não fazem mal; a menos, podem dar resultado errado.

\subsection{Nomes, cálculos e nulos}

\begin{lstlisting}
select titulo, preco * 0.9 as preco_desconto from livro
select titulo, preco * 0.9 "preço com desconto" from livro
select titulo from livro where preco is null
\end{lstlisting}

O \texttt{as} é opcional, e as aspas só são necessárias quando o nome tem
espaço. O apelido é útil sobretudo em colunas calculadas, que sairiam sem
nome. No SQL Server, o apelido não vale no \texttt{where} (é preciso
repetir a expressão), mas vale no \texttt{order by}.

Para testar nulo, use \texttt{is null} e \texttt{is not null}: a
comparação \texttt{preco = null} não traz linha nenhuma.

\subsection*{Complemento: a lógica de três valores}

\textcite[seção 5.1.1]{Elmasri2011} explica por que \texttt{= null} não
funciona. Um nulo pode significar valor desconhecido, valor omitido de
propósito ou atributo que não se aplica, e a SQL não distingue os três
casos. Toda comparação com nulo dá \textsc{desconhecido}, e não
verdadeiro ou falso. A SQL usa então uma lógica de três valores:

\begin{itemize}
	\item \textsc{falso} \texttt{and} qualquer coisa dá \textsc{falso};
	\textsc{verdadeiro} \texttt{and} \textsc{desconhecido} dá
	\textsc{desconhecido};
	\item \textsc{verdadeiro} \texttt{or} qualquer coisa dá
	\textsc{verdadeiro}; \textsc{falso} \texttt{or} \textsc{desconhecido}
	dá \textsc{desconhecido};
	\item \texttt{not} \textsc{desconhecido} continua
	\textsc{desconhecido}.
\end{itemize}

O \texttt{where} só devolve as linhas em que a condição é
\textsc{verdadeiro}. Por isso \texttt{not in} tem uma armadilha: se a
lista (ou a subconsulta) contém um nulo, a condição nunca é verdadeira e
a consulta volta vazia.

\subsection{Ordenação}

\begin{lstlisting}
select genero, titulo, preco from livro order by genero, preco desc
\end{lstlisting}

O padrão é \texttt{asc}. No SQL Server, os nulos vêm primeiro na ordem
ascendente e por último na descendente. Ordenar pela segunda coluna só
muda algo se a primeira tiver valores repetidos; ordenar por chave e
depois por outra coluna gasta recurso à toa.

\subsection{Funções de agregação, \texttt{group by} e \texttt{having}}

\begin{lstlisting}
select count(*) from livro              -- linhas
select count(preco) from livro          -- linhas com preço não nulo
select avg(preco), min(preco), max(preco), sum(preco) from livro
select avg(distinct preco) from livro

select genero, avg(preco) as media
from livro
group by genero
having avg(preco) > 40
\end{lstlisting}

\begin{itemize}
	\item \texttt{count(*)} conta linhas; \texttt{count(coluna)} e as
	outras funções ignoram nulos.
	\item Sem \texttt{group by}, \texttt{select genero, avg(preco)} dá
	erro: a média devolve um valor e \texttt{genero}, vários. Toda coluna
	fora de função de agregação precisa estar no \texttt{group by}.
	\item Agrupar por chave não faz sentido: cada grupo teria uma linha só.
	\item O \texttt{where} filtra linhas \emph{antes} do agrupamento e não
	enxerga os grupos; para filtrar grupos existe o \texttt{having}. A
	condição do \texttt{having} pode usar uma função diferente da que está
	no \texttt{select}.
\end{itemize}

\section{Junções}

O \textbf{produto cartesiano} (\texttt{cross join}) combina cada linha
de uma tabela com todas as da outra: 18 títulos e 8 editoras dão 144
linhas. Sozinho, quase nunca tem uso. A \textbf{junção}\index{junção}
é o produto seguido de uma condição, geralmente igualdade entre chave
estrangeira e chave primária. O algoritmo, porém, não monta o produto
inteiro para filtrar depois: testa a condição a cada combinação e
descarta na hora o que não serve.

\begin{lstlisting}[caption={Junção interna com duas e três tabelas.}]
select l.titulo, a.nome
from livro l inner join autor a
     on l.codautor = a.codautor
where l.preco > 50

select le.nome, l.titulo, e.data
from leitor le inner join emprestimo e
     on le.codleitor = e.codleitor
     inner join livro l
     on e.codlivro = l.codlivro
\end{lstlisting}

\begin{itemize}
	\item A condição de junção fica no \texttt{on}; os filtros ficam no
	\texttt{where}.
	\item Junta-se de duas em duas tabelas, seguindo os relacionamentos.
	Uma tabela já incluída fica visível para os \texttt{on} seguintes.
	\item Uma tabela não tem colunas com o mesmo nome, mas o resultado de
	uma junção pode ter. Sem o prefixo (\texttt{l.codautor}), dá erro de
	ambiguidade. O apelido da tabela deve ser curto.
	\item \texttt{join} sozinho é \texttt{inner join}.
	\item A sintaxe antiga, com vírgula no \texttt{from} e a condição no
	\texttt{where}, ainda funciona, mas mistura junção com filtro e fica
	confusa com três tabelas ou mais. A professora pede a notação
	\texttt{inner join ... on}.
\end{itemize}

\subsection{Junções externas}

\begin{lstlisting}
-- todos os leitores, inclusive os sem empréstimo
select le.nome, e.codemprestimo
from leitor le left join emprestimo e
     on le.codleitor = e.codleitor

-- só os leitores sem empréstimo
select le.nome
from leitor le left join emprestimo e
     on le.codleitor = e.codleitor
where e.codemprestimo is null
\end{lstlisting}

\begin{figure}[h]
\centering
\begin{tikzpicture}[
	tab/.style={draw=bdcinza, font=\scriptsize, minimum width=17mm, minimum height=5mm, inner sep=1pt},
	ok/.style={tab, fill=bdazulclaro},
	nul/.style={tab, fill=white, text=bdcinza}]
	\node[font=\scriptsize\bfseries] at (0,0.6) {leitor};
	\node[tab] (a1) at (0,0) {1 Ana};
	\node[tab, below=0pt of a1] (a2) {2 Bruno};
	\node[tab, below=0pt of a2] (a3) {3 Carla};
	\node[font=\scriptsize\bfseries] at (2.3,0.6) {empréstimo};
	\node[tab] (b1) at (2.3,0) {100 \textrightarrow{} 1};
	\node[tab, below=0pt of b1] (b2) {101 \textrightarrow{} 3};
	\node[tab, below=0pt of b2] (b3) {102 \textrightarrow{} 9};
	% inner
	\node[font=\scriptsize\bfseries] at (5.2,0.6) {inner join};
	\node[ok] (i1) at (5.2,0) {Ana | 100};
	\node[ok, below=0pt of i1] (i2) {Carla | 101};
	% left
	\node[font=\scriptsize\bfseries] at (7.5,0.6) {left join};
	\node[ok] (l1) at (7.5,0) {Ana | 100};
	\node[nul, below=0pt of l1] (l2) {Bruno | null};
	\node[ok, below=0pt of l2] (l3) {Carla | 101};
\end{tikzpicture}
\caption{Junção interna e externa à esquerda. Bruno não tem empréstimo e
só aparece no \texttt{left join}, com nulo. O empréstimo 102 aponta para
um leitor inexistente, situação que a chave estrangeira impediria; ele só
apareceria num \texttt{right} ou \texttt{full join}.}
\end{figure}

O \texttt{left join} mantém todas as linhas da tabela da esquerda e
preenche com nulos o que não tem correspondência; o \texttt{right join}
faz o mesmo com a da direita, e o \texttt{full join} com as duas. A
palavra \texttt{outer} é opcional. Para achar as linhas sem
correspondência, testa-se \texttt{is null} numa coluna da tabela do
outro lado, que nem precisa estar no \texttt{select}.

\subsection{O erro que conta na nota}

O alerta com mais peso de avaliação da aula: \emph{não juntar tabela
desnecessária}. Junção é cara. Se o código que já está na tabela de
relacionamento basta, a consulta para ali. Numa consulta otimizada,
juntar uma tabela que não contribui com nenhuma coluna nem filtro conta
como erro, mesmo que o resultado saia certo.

\begin{lstlisting}
-- desnecessário: autor não contribui com nada
select l.titulo, l.codautor
from livro l inner join autor a on l.codautor = a.codautor

-- suficiente
select titulo, codautor from livro
\end{lstlisting}

\section{Subconsultas}

\begin{lstlisting}
-- livros da editora "Nova Lua", sem junção
select titulo from livro
where codeditora = (select codeditora from editora
                    where nome = 'Nova Lua')

-- livros de autores com sobrenome Ramos (pode haver vários)
select titulo from livro
where codautor in (select codautor from autor
                   where sobrenome = 'Ramos')
\end{lstlisting}

A posição da professora é clara:

\begin{itemize}
	\item a subconsulta que vale a pena é a do \texttt{where} que
	\emph{substitui} uma junção: são dois passos simples, sem produto
	cartesiano;
	\item use \texttt{=} só quando tiver certeza de que a subconsulta
	devolve um único valor; senão, \texttt{in};
	\item a consulta externa não enxerga colunas das tabelas internas. Se
	essas colunas são necessárias, a solução é a junção;
	\item subconsultas no \texttt{select} ou no \texttt{from} precisam ser
	muito bem pensadas. Na opinião dela, perdem para a junção na maioria
	dos casos, porque cada uma pode carregar sua própria junção. Ela
	disse ver esse uso em toda parte, possivelmente em código gerado por
	IA, e pediu senso crítico. O tema volta na disciplina de otimização.
\end{itemize}

\textcite[seção 5.1]{Elmasri2011} vai além da aula com subconsultas
correlacionadas, \texttt{exists} e \texttt{unique}.

\section{Transações}

Uma \textbf{transação}\index{transação} é um conjunto de operações que o
usuário enxerga como uma só: ou tudo é feito, ou nada. Exemplos da aula:
uma transferência entre contas (débito e crédito); uma venda com baixa
de estoque; o cadastro de uma pessoa física, que grava a parte genérica
e a específica.

\subsection{Estados}

\begin{description}
	\item[Ativa:] desde o \texttt{begin}, executando leituras e escritas.
	\item[Parcialmente efetivada:] depois da última operação, mas antes do
	\texttt{commit}. Ainda pode falhar (uma queda de energia, por
	exemplo).
	\item[Efetivada:] depois do \texttt{commit}.
	\item[Falha:] depois de um erro, com \texttt{rollback} (ou
	\emph{abort}, sinônimos).
	\item[Encerrada:] o destino de todas, com sucesso ou com falha.
\end{description}

\begin{figure}[h]
\centering
\begin{tikzpicture}[estado/.style={caixa, rounded corners=8pt, minimum width=24mm}]
	\node[estado] (at) at (0,0) {Ativa};
	\node[estado] (pe) at (5,0) {Parcialmente\\efetivada};
	\node[estado] (ef) at (5,-2) {Efetivada};
	\node[estado] (fa) at (0,-2) {Falha};
	\node[estado, fill=bdroxo, text=white] (en) at (2.5,-3.6) {Encerrada};
	\draw[seta] (-2.6,0) -- node[rotulo, above] {begin} (at);
	\draw[seta] (at) -- node[rotulo, above] {última operação} (pe);
	\draw[seta] (pe) -- node[rotulo, right] {commit} (ef);
	\draw[seta] (at) -- node[rotulo, left] {erro} (fa);
	\draw[seta] (pe) -- node[rotulo, above, sloped] {falha} (fa);
	\draw[seta] (ef) -- (en);
	\draw[seta] (fa) -- node[rotulo, left] {rollback} (en);
	\draw[seta] (at) edge[loop above, looseness=5] node[rotulo, above] {leituras e escritas} (at);
\end{tikzpicture}
\caption{Estados de uma transação. Com sucesso ou com falha, toda
transação precisa chegar a encerrada, que libera os bloqueios.}
\end{figure}

A professora insistiu que toda transação precisa chegar ao estado
encerrada: uma transação aberta mantém os bloqueios e deixa os outros
usuários esperando. O \texttt{rollback} ``dá a impressão de desfazer'',
mas é sobretudo o aviso de encerramento.

\subsection{Propriedades ACID}

\begin{description}
	\item[Atomicidade:] tudo ou nada.
	\item[Consistência:] a transação leva o banco de um estado consistente
	a outro.
	\item[Isolamento:] cada transação se comporta como se executasse
	sozinha.
	\item[Durabilidade:] o que foi efetivado não se perde, nem com falha.
\end{description}

Na divisão da aula, atomicidade e consistência dependem de quem escreve
a transação: o SGBD obedece, mas só o projetista sabe o que precisa ser
agrupado. Se o débito de uma transferência for esquecido, o SGBD executa
o crédito sem reclamar. Isolamento e durabilidade são responsabilidade
do SGBD, que enxerga tudo o que está rodando.

\textcite[seção 21.3]{Elmasri2011} divide um pouco diferente: o
subsistema de recuperação garante atomicidade e durabilidade, o controle
de concorrência garante isolamento, e a consistência cabe aos
programadores e ao módulo de restrições de integridade. As duas visões
se completam: o SGBD garante o ``tudo ou nada'' daquilo que o
programador colocou dentro da transação.

\subsection{Conflito e bloqueio}

Duas operações \textbf{conflitam} quando pertencem a transações
diferentes, acessam o mesmo item (tabela, página ou linha, conforme a
granularidade) e pelo menos uma delas é escrita. Leituras não conflitam
entre si. Diante de um conflito, o SGBD serializa as operações.

O mecanismo usual é o \textbf{bloqueio}\index{bloqueio}: compartilhado
para leitura e exclusivo para escrita. Bloquear o banco inteiro seria
absurdo; a tabela é fácil de bloquear; a linha gera menos conflitos, mas
custa mais para gerenciar. A página é o meio-termo, e é a granularidade
mais comum. Bloqueios explícitos, escritos pelo programador, existem,
mas o desejável é deixar o SGBD bloquear sozinho.

Um caso comentado em aula: uma escrita precisa de bloqueio exclusivo e
pode esperar para sempre enquanto chegam leituras com bloqueio
compartilhado (\emph{livelock}). A solução é que, quando há uma escrita
na fila, nenhuma nova leitura recebe bloqueio compartilhado.

Regra prática da professora: transação enxuta, só com o que precisa ser
agrupado, e com fim claro em \texttt{commit} ou \texttt{rollback}.

\subsection{Transações no SQL Server}
\labsec{cet615-transacoes-sqlserver}

\begin{itemize}
	\item O modo padrão é \emph{autocommit}: cada comando é uma transação
	e é efetivado sozinho se não der erro.
	\item A transação explícita começa com \texttt{begin tran} (ou
	\texttt{begin transaction}) e termina com \texttt{commit} ou
	\texttt{rollback}. Não existe \texttt{end transaction}.
	\item No modo implícito, cada \texttt{commit} ou \texttt{rollback} já
	abre a próxima transação. A professora desaconselha: exige muito
	domínio e tende a segurar bloqueios por mais tempo.
	\item \texttt{save tran} cria um ponto de retorno parcial. Ela não
	vê sentido nele numa transação bem projetada.
	\item O bloqueio é implícito; a granularidade padrão é a página.
	\item \texttt{@@rowcount} é a variável do sistema com o número de
	linhas afetadas pelo último comando. Ela usa \texttt{> 0} em vez de
	\texttt{= 1} para valer também em \texttt{update} de várias linhas.
	\item Blocos são \texttt{begin ... end}. Não existe \texttt{end if}.
\end{itemize}

A aula não usa \texttt{try...catch}: a lógica é escrita com
\texttt{if}, que todo SGBD tem, para deixar visível o caminho de cada
falha. O padrão de uma transferência:

\begin{lstlisting}[caption={Transferência entre contas com \texttt{@@rowcount}.}]
begin transaction
update conta set saldo = saldo - 100 where numero = 1
if @@rowcount > 0               -- débito feito
begin
   update conta set saldo = saldo + 100 where numero = 2
   if @@rowcount > 0            -- crédito feito
      commit transaction
   else
      rollback transaction
end
else
   rollback transaction
\end{lstlisting}

Todo caminho termina em \texttt{commit} ou \texttt{rollback}, inclusive
o primeiro \texttt{else}: mesmo sem nada gravado, a transação foi aberta
e precisa ser encerrada. Com chaves estrangeiras, a ordem importa:
insere-se primeiro a linha referenciada (cliente antes de pessoa física,
conta antes de cliente--conta). Na demonstração da aula, a segunda
inserção falhou e a primeira não ficou no banco. Sem a transação, ela
teria ficado gravada sozinha.

\subsection*{Complemento: o que o \texttt{rollback} faz por dentro}

Em aula, a professora explicou que as escritas ficam no log e só vão
para o banco no \texttt{commit}; o \texttt{rollback} apenas as descarta.
É a estratégia de atualização adiada, que \textcite[cap. 23]{Elmasri2011}
chama de NO-UNDO/REDO. O livro descreve também a atualização imediata,
em que páginas alteradas podem ir para o disco antes do \texttt{commit}
e o \texttt{rollback} precisa desfazê-las.

O SQL Server segue a segunda linha\sidecite{MSTransactionLog}. Ele usa
\emph{write-ahead logging}: nenhuma página alterada vai para o disco
antes do registro de log correspondente. As páginas são alteradas na
memória (\emph{buffer cache}) e gravadas depois, por exemplo num
\emph{checkpoint}, que pode acontecer antes do \texttt{commit}. Por isso
o log guarda a imagem anterior de cada alteração, e o \texttt{rollback}
a reaplica. Na recuperação depois de uma queda, o SQL Server refaz o
que foi efetivado e não chegou ao disco e desfaz o que ficou sem
\texttt{commit}. O resultado visto pelo usuário é o mesmo da explicação
da aula: o que não foi efetivado não fica no banco.

A documentação reforça também o alerta sobre transações abertas: o log
não pode ser truncado antes da transação ativa mais antiga, então uma
transação esquecida faz o log crescer sem parar.
